\subsection{Differentialgleichung einer RC-Reihenschaltung}
\label{subsec:differentialgleichung-rc-schaltung}

Wir betrachten eine RC-Reihenschaltung mit einer konstanten Eingangsspannung
$U_0$, einem Widerstand $R$ und einem Kondensator $C$.
Die Spannung am Widerstand wird mit $U_R(t)$ bezeichnet, die Spannung am
Kondensator mit $U_C(t)$. Da Widerstand und Kondensator in Reihe geschaltet
sind, fließt durch beide Bauelemente derselbe Strom $i(t)$.
\FloatBarrier


\begin{center}
	\begin{tikzpicture}[scale=1.0, european resistors]
		
		% Coordinates
		\coordinate (A) at (0,0);
		\coordinate (B) at (3.5,0);
		\coordinate (C) at (5,0);
		\coordinate (D) at (5,-2);
		\coordinate (E) at (0,-2);
		
		% Source
	
		% Gleichspannungsquelle
	
	\draw (A) to[battery1] (E);
		\draw[<-,blue] (-0.6,-0.4) -- (-0.6,-1.4);
	\node[above,blue] at (-1,-1.25) {$U_0(t)$};
		% Resistor
		\draw (A) -- (2.2,0);
		\draw (2.2,0) to[R, l={$R$}] (3.5,0);
		
		% Capacitor
		\draw (3.5,0) -- (C);
		\draw (C) to[C, l={$C$}] (D);
	     \node[right] at (4.5,-0.6) {$+$};
		% Ground return
		\draw (D) -- (E);
		
		% Capacitor voltage arrow
		\draw[->,green] (5.9,-0.4) -- (5.9,-1.4);
		\node[right,green] at (6.0,-1.0) {$U_C(t)$};
		
		% Rissistant arrow
		\draw[->,blue] (2.2,0.8) -- (3.2,0.8);
		\node[above,blue] at (2.8,0.8) {$U_R(t))$};
			% Current arrow
		\draw[->,red] (0.2,0.2) -- (1.2,0.2);
		\node[above,red] at (0.8,0.3) {$i(t))$};
		
	\end{tikzpicture}
\end{center}




Nach der Maschenregel gilt:

\[
U_0 = U_R(t) + U_C(t)
\]

Für den Widerstand gilt nach dem Ohmschen Gesetz:

\[
U_R(t) = R\,i(t)
\]

Für den Kondensator gilt:

\[
i(t) = C\,\frac{dU_C(t)}{dt}
\]

Setzt man diesen Ausdruck für den Strom in die Gleichung des Widerstands ein,
so erhält man:

\[
U_R(t) = R C\,\frac{dU_C(t)}{dt}
\]

Einsetzen in die Maschengleichung liefert:

\[
U_0 = R C\,\frac{dU_C(t)}{dt} + U_C(t)
\]

Damit lautet die Differentialgleichung der RC-Reihenschaltung:

\[
R C\,\frac{dU_C(t)}{dt} + U_C(t) = U_0
\]

Da das Produkt aus Widerstand und Kapazität die Zeitkonstante der Schaltung ist,

\[
\tau = R C,
\]

kann man die Gleichung auch schreiben als:

\[
\tau\,\frac{dU_C(t)}{dt} + U_C(t) = U_0
\]

oder in normierter Form:

\[
\frac{dU_C(t)}{dt} + \frac{1}{\tau}U_C(t)
=
\frac{U_0}{\tau}
\]

Es handelt sich also um eine lineare Differentialgleichung erster Ordnung
mit konstanter rechter Seite.

\subsubsection{Allgemeine Lösung}
\label{subsubsec:allgemeine-loesung-rc-schaltung}

Die allgemeine Lösung dieser Differentialgleichung besitzt die Form

\[
U_C(t) = K_1 + K_2 e^{-\frac{t-t_0}{R C}}
\]

oder mit der Zeitkonstanten $\tau = RC$:

\[
U_C(t) = K_1 + K_2 e^{-\frac{t-t_0}{\tau}}
\]

Da $U_0$ konstant ist, ist die stationäre Lösung

\[
K_1 = U_0.
\]

Damit ergibt sich:

\[
U_C(t) = U_0 + K_2 e^{-\frac{t-t_0}{R C}}
\]

bzw.

\[
U_C(t) = U_0 + K_2 e^{-\frac{t-t_0}{\tau}}.
\]

Die Konstante $K_2$ wird durch die Anfangsbedingung bestimmt.

Sei zum Zeitpunkt $t=t_0$ die Kondensatorspannung

\[
U_C(t_0)=U_{C0}.
\]

Dann folgt aus der allgemeinen Lösung:

\[
U_{C0}=U_0+K_2.
\]

Also ist

\[
K_2 = U_{C0}-U_0.
\]

Damit erhält man die Lösung für beliebige Anfangsspannung des Kondensators:

\[
\boxed{
	U_C(t)=U_0+\left(U_{C0}-U_0\right)e^{-\frac{t-t_0}{R C}}
}
\]

oder mit $\tau=RC$:

\[
\boxed{
	U_C(t)=U_0+\left(U_{C0}-U_0\right)e^{-\frac{t-t_0}{\tau}}
}
\]

\subsubsection{Spezialfall: ungeladener Kondensator}
\label{subsubsec:ungeladener-kondensator-rc-schaltung}

Ist der Kondensator zu Beginn ungeladen, so gilt:

\[
U_C(t_0)=0.
\]

Damit ist

\[
U_{C0}=0.
\]

Einsetzen in die allgemeine Lösung ergibt:

\[
U_C(t)
=
U_0+\left(0-U_0\right)e^{-\frac{t-t_0}{R C}}.
\]

Also folgt:

\[
\boxed{
	U_C(t)=U_0\left(1-e^{-\frac{t-t_0}{R C}}\right)
}
\]

oder mit der Zeitkonstanten $\tau$:

\[
\boxed{
	U_C(t)=U_0\left(1-e^{-\frac{t-t_0}{\tau}}\right)
}
\]

Dies ist die klassische Aufladekurve des Kondensators.

\subsubsection{Spannung am Widerstand}
\label{subsubsec:spannung-am-widerstand-rc-schaltung}

Aus der Maschengleichung

\[
U_0 = U_R(t)+U_C(t)
\]

folgt unmittelbar:

\[
U_R(t)=U_0-U_C(t).
\]

Für den ungeladenen Kondensator ergibt sich daher:

\[
U_R(t)
=
U_0
-
U_0\left(1-e^{-\frac{t-t_0}{R C}}\right).
\]

Damit erhält man:

\[
\boxed{
	U_R(t)=U_0 e^{-\frac{t-t_0}{R C}}
}
\]

oder mit $\tau=RC$:

\[
\boxed{
	U_R(t)=U_0 e^{-\frac{t-t_0}{\tau}}
}
\]

Die Widerstandsspannung fällt also exponentiell ab, während die
Kondensatorspannung exponentiell gegen $U_0$ ansteigt.

\subsubsection{Bedeutung der Zeitkonstante}
\label{subsubsec:bedeutung-zeitkonstante-rc-schaltung}

Die Größe

\[
\boxed{
	\tau = R C
}
\]

heißt Zeitkonstante der RC-Schaltung.

Nach einer Zeitspanne

\[
t-t_0=\tau
\]

hat der Kondensator die Spannung

\[
U_C(t_0+\tau)
=
U_0\left(1-e^{-1}\right)
\]

erreicht. Da

\[
1-e^{-1}\approx 0{,}632
\]

gilt:

\[
U_C(t_0+\tau) \approx 0{,}632\,U_0.
\]

Nach einer Zeitkonstante hat der Kondensator also etwa $63{,}2\,\%$ seiner
Endspannung erreicht.

Umgekehrt ist die Widerstandsspannung nach einer Zeitkonstante auf

\[
U_R(t_0+\tau)=U_0 e^{-1}
\]

abgefallen. Da

\[
e^{-1}\approx 0{,}368
\]

gilt:

\[
U_R(t_0+\tau)\approx 0{,}368\,U_0.
\]

Nach einer Zeitkonstante beträgt die Spannung am Widerstand also nur noch etwa
$36{,}8\,\%$ ihres Anfangswerts.

\begin{DidacticBox}[Merksatz zur RC-Schaltung]
	\label{box:merksatz-rc-schaltung}
	Bei einer RC-Reihenschaltung mit konstanter Eingangsspannung $U_0$ steigt die
	Kondensatorspannung exponentiell an, während die Widerstandsspannung
	exponentiell abfällt:
	
	\[
	U_C(t)=U_0\left(1-e^{-\frac{t-t_0}{\tau}}\right),
	\qquad
	U_R(t)=U_0 e^{-\frac{t-t_0}{\tau}}.
	\]
	
	Die Zeitkonstante
	
	\[
	\tau=RC
	\]
	
	bestimmt, wie schnell dieser Vorgang abläuft.
\end{DidacticBox}